
\section{Relations Between Cartan Subgroups and Cartan Subalgebras}
\SGARefTarget{sga3:xiv:5}

Applying \SGARefLink{sga3:xiv:4.12}{4.12} to the adjoint representation of \(G\), we obtain:

\paragraph{\textbf{Theorem 5.1.}}
\SGARefTarget{sga3:xiv:5.1}
{\itshape Let \(G\) be a smooth algebraic group over an infinite
field, let \(T\) be a maximal torus of \(G\), and let
\(\mathfrak g\supset\mathfrak t\) be their Lie algebras. Then
\(\mathfrak t\) contains a regular element of \(\mathfrak g\).}

\paragraph{\textbf{Corollary 5.2.}}
\SGARefTarget{sga3:xiv:5.2}
{\itshape Let \(G\) be a smooth \(S\)-prescheme in groups of finite
presentation, and let \(\mathfrak g\) be its Lie algebra.

\smallskip
\noindent
\textup{a)} Conditions \((\mathrm C_0)\) through \((\mathrm C_2)\)
of \SGARefLink{sga3:xiv:2.9}{2.9} for \(\mathfrak g\) are equivalent. In particular, if the
infinitesimal rank of the fibers of \(\mathfrak g\) at the points
of \(S\) is locally constant, then, étale locally,
\(\mathfrak g\) admits a Cartan subalgebra and hence, by
\SGARefLink{sga3:xiv:3.9}{3.9}\textup{(a)}, \(G\) admits a subgroup of type \((\mathrm C)\).

\smallskip
\noindent
\textup{b)} Let \(H\) be a subprescheme in groups of \(G\), smooth
and of finite presentation over \(S\), whose fibers are connected
and at every \(s\in S\) have the same reductive rank as \(G\);
for example, \(H\) may be a maximal torus or a Cartan subgroup of
\(G\). Let \(D\) be a subgroup of type \((\mathrm C)\) of \(G\).
Then
\[
  H\subset D\quad\Longleftrightarrow\quad
  \mathfrak h\subset\mathfrak d.
\]

\smallskip
\noindent
\textup{c)} Suppose that condition \((\mathrm C_0)\) holds, that is,
the infinitesimal rank of \(G\) is locally constant. Let \(H\) be a
subprescheme in groups of \(G\), smooth and of finite presentation
over \(S\), with connected nilpotent fibers, and having at every
point the same reductive rank as \(G\); for example, \(H\) may be a
maximal torus or a Cartan subgroup of \(G\). Then \(H\) is contained
in a unique subgroup \(D\) of type \((\mathrm C)\) of \(G\).

\smallskip
\noindent
\textup{d)} If \(G\) admits a Cartan subgroup fpqc locally, then so
does every subgroup \(D\) of type \((\mathrm C)\) of \(G\).}

\smallskip
\noindent
\emph{Proof.}
\textup{a)}
Assume \((\mathrm C_0)\), and let us prove \((\mathrm C_2)\), that
is, that every quasi-regular section \(a\) of \(\mathfrak g\) is
regular. As usual, we reduce first to \(S\) affine noetherian and
then, since the question is infinitesimal by \SGARefLink{sga3:xiv:2.15}{2.15}\textup{(b)}, to
\(S\) local Artinian. In this last case \((\mathrm C_0)\) is
automatic. We may moreover suppose that the residue field \(k\)
of \(S\) is infinite.

By \SGARefLink{sga3:xiv:3.7}{3.7}, it is enough to show that \(\mathfrak g\) admits a Cartan
subalgebra. Let \(T\) be a maximal torus of \(G\). By \SGARefLink{sga3:xiv:5.1}{5.1} there is
a quasi-regular element \(t\) in the Lie algebra \(\mathfrak t\)
of \(T\). We shall show that \(t\) is regular.

Consider the linear representation of \(T\) on \(\mathfrak g\)
induced by the adjoint representation of \(G\). There is a finite
set of characters \((u_i)_{i\in I}\) of \(T\) and a direct-sum
decomposition
\[
  \mathfrak g=\bigoplus_{i\in I}\mathfrak g_i
\]
into \(T\)-stable submodules such that \(T\) acts on
\(\mathfrak g_i\) through \(u_i\)
\textup{(I, \S4.7.3)}. Let
\[
  u_i'\colon\mathfrak t\longrightarrow A
\]
be the homomorphism of Lie algebras induced by
\(u_i\colon T\to\mathbb G_m\), where \(A\) is the ring of \(S\).
Let \(u_{i0}\) and \(u_{i0}'\) be their reductions under
\(A\to k\). Put
\[
  I'=\{i\in I\mid u_{i0}'\ne0\},
  \qquad I''=I\setminus I'.
\]
The regularity of \(t_0\) is expressed by
\[
  u_{i0}'(t_0)\ne0\qquad(i\in I'),
\]
and therefore \(u_i'(t)\) is invertible for \(i\in I'\).

The Cartan subalgebra of \(\mathfrak g_0\) defined by \(t_0\), that
is, the kernel of the semisimple endomorphism
\(\operatorname{ad}(t_0)\), is
\[
  \bigoplus_{i\in I''}(\mathfrak g_i)_0.
\]
Set
\[
  \mathfrak d=\bigoplus_{i\in I''}\mathfrak g_i.
\]
Then \(\operatorname{ad}(t)\) is nilpotent on \(\mathfrak d\), and
is an automorphism of
\[
  \mathfrak g/\mathfrak d
  \simeq\bigoplus_{i\in I'}\mathfrak g_i.
\]
By \SGARefLink{sga3:xiv:2.6}{2.6}, \(t\) is regular.

\smallskip
\noindent
\textup{b)}
This follows from \SGARefLink{sga3:xiv:3.12}{3.12}\textup{(a)}. By \SGARefLink{sga3:xiv:5.1}{5.1}, \(H\) satisfies the
required hypothesis that every geometric fiber of \(\mathfrak h\)
contains a regular element of the corresponding fiber of
\(\mathfrak g\).

\smallskip
\noindent
\textup{c)}
By \textup{b)}, it is enough to show that \(\mathfrak h\) is
contained in a unique Cartan subalgebra of \(\mathfrak g\). By
\textup{a)}, \(\mathfrak g\) satisfies \((\mathrm C_2)\). We are
therefore reduced to the following lemma.

\paragraph{\textbf{Lemma 5.3.}}
\SGARefTarget{sga3:xiv:5.3}
{\itshape Let \(\mathfrak g\) be a Lie algebra over \(S\), locally
free of finite type as a module, and satisfying condition
\((\mathrm C_2)\) of \SGARefLink{sga3:xiv:2.9}{2.9}. Let \(\mathfrak h\) be a Lie subalgebra
of \(\mathfrak g\) such that \(\mathfrak h\) is locally a direct
summand, is locally nilpotent, and, for every \(s\in S\), the
geometric fiber \(\mathfrak h_{\bar s}\) contains a regular element
of \(\mathfrak g_{\bar s}\). Then \(\mathfrak h\) is contained in
a unique Cartan subalgebra of \(\mathfrak g\).}

In the situation at hand, \(\mathfrak h\) has these properties: it
is locally nilpotent because \(H\) has nilpotent fibers, and the
condition on regular elements follows from \SGARefLink{sga3:xiv:5.1}{5.1}.

Since \SGARefLink{sga3:xiv:5.3}{5.3} is fpqc local, it is enough, at a point \(s\in S\) with
\(\kappa(s)\) infinite, to find an open neighborhood \(U\) of \(s\)
on which existence and uniqueness hold after every base change
\(S'\to S\) factoring through \(U\). Choose a regular element of
\(\mathfrak g\otimes\kappa(s)\) lying in
\(\mathfrak h\otimes\kappa(s)\), and extend it to a section \(a\)
of \(\mathfrak h\) over an open neighborhood \(U\) of \(s\).
Condition \((\mathrm C_2)\) and \SGARefLink{sga3:xiv:2.10}{2.10} allow us to shrink \(U\) so
that \(a\) is regular. We may thus suppose \(U=S\).

Every Cartan subalgebra of \(\mathfrak g\) containing
\(\mathfrak h\) contains \(a\), and hence equals
\[
  \mathfrak d=\Nil(a,\mathfrak g)
\]
by \SGARefLink{sga3:xiv:2.6}{2.6}; this proves uniqueness. Conversely,
\(\mathfrak h\subset\mathfrak d\) because \(\mathfrak h\) is
locally nilpotent, which proves existence.

\smallskip
\noindent
\textup{d)}
This is a special case of \SGARefLink{sga3:xii:7.9}{XII~7.9}\textup{\SGARefLink{sga3:xii:7.9:item:d:02}{(d)}}.
